\section{\sys Design}
\label{sec:design}

In this section, we present \sys, the first decentralized storage system that simultaneously achieves the three properties in \cref{sec:approach:model}.
The main novelty of \sys is a distributed storage design that carefully applies random sampling (\cref{sec:approach:case}) to both data placement and erasure coding.
The design assumes a BFT blockchain that orders and executes smart contracts (\ie, transactions) with linearizability~\cite{linearizability} guarantee;
\sys is agnostic to the protocol details of the blockchain.
We also assume the blockchain provides an incentive token that is economically favorable to rational nodes.

\subsection{VRF and Rateless EC Primer}
\label{sec:design:primer}

\sys applies two algorithms in its design: verifiable random function~\cite{vrf} and verifiable rateless erasure code~\cite{verify-rateless-ec}.
We first give a primer of the two algorithms.

\paragraph{\textbf{Verifiable random function (VRF)}}
A VRF~\cite{vrf} exposes a $\VRFGen$ and a $\VRFVerify$ function.
$\VRFGen$ takes a secret key $\sk$ and an input $x$, and outputs a pseudorandom value $y$ and a proof $\pi$.
$\VRFVerify$ takes a public key $\pk$, an input $x$, a value $y$, and a proof $\pi$, and outputs a boolean indicating whether $y$ is the correct $\VRFGen$ output for $x$ under $\pk$.

A VRF satisfies the following properties:
\begin{itemize}
    \item \textbf{Pseudorandomness:} For any input $x$, the output $y$ is indistinguishable from random to any efficient adversary that does not know the secret key $\sk$.
    \item \textbf{Uniqueness:} For any input $x$, there is a unique output $y$ that can be generated by $\VRFGen$ with the corresponding proof $\pi$.
    \item \textbf{Verifiability:} Given the public key $\pk$, input $x$, output $y$, and proof $\pi$, anyone can efficiently verify that $y$ is the correct VRF output for $x$ under $\pk$ using $\VRFVerify$.
\end{itemize}

\paragraph{\textbf{Verifiable rateless erasure code}}
Rateless codes are a family of erasure codes that has a special property.
A rateless code generates an \emph{infinite} sequence of encoding symbols from $k$ source blocks.
Any $k + \epsilon$ encoding symbols can be used to reconstruct the original data.
This differs from traditional maximum distance separable (MDS) codes, such as Reed-Solomon codes, which produce a fixed number of $n$ encoding symbols from the $k$ source blocks.

Internally, a rateless code such as LT codes~\cite{lt_code} and raptor codes~\cite{raptor_code} split the input data into blocks.
For each encoding symbol, it performs an XOR on some subset of the blocks.
With sufficient number of blocks, the number of possible encoding symbols (exponential to the number of blocks) is practically infinite.
To decode, the decoder performs XOR on the encoding symbols to generate symbols with lesser degrees (degree is the number of original blocks XORed to produce the symbol).
The decoder repeats the process until it generates all degree one blocks, \ie, the original blocks.

Using traditional rateless codes, a decoder cannot verify the integrity of a received encoding symbol;
performing decoding using a tampered symbol would lead to integrity violation of the data.
Verifiable rateless code~\cite{verify-rateless-ec} enables any third party to validate the integrity of a received symbol.
The scheme is also space efficient.
The third party only needs to possess a small hash of the original data when performing the validation.

\lijl{Like VRF, can show the verifiable rateless EC interface here, so that later text can refer to it.}

\subsection{\sys Protocol}
\label{sec:design:proto}

An \sys deployment consists of three main components: storage nodes that provide permanent storage, clients which issue \textsc{Store} and \textsc{Retrieve} requests, and \sys smart contracts deployed on the blockchain.
\cref{alg:contract}, \cref{alg:client}, and \cref{alg:node} show the state and the pseudocode for the smart contracts, the clients, and the storage nodes, respectively.
We only require the storage nodes to submit transactions and learn the finalized transaction blocks on the blockchain, that is, storage nodes and blockchain full nodes can be disjoint.

\subsubsection{Node Management}
\label{sec:design:proto:mgmt}
\sys runs in a permissionless setting.
Storage nodes can join or leave the system at any time.
To join the system, a storage node submits a \textsc{Stake} transaction with its public key and staking tokens, as shown in \cref{alg:contract}.
The on-chain state maintains the set of currently staking nodes, which are the active storage nodes, in the system.
The transaction adds the node to the set, and locks up the staked tokens.
If the node later fails to pass the storage verification tests~\cref{sec:design:proto:store}, the staked tokens will be slashed.
To leave the system, an existing node submits a \textsc{UnStake} transaction, which removes the node from the staking set and returns the locked tokens.
\lijl{Can add the unstake transaction to the pseudocode.}
By checking the $\st{stakingSet}$ on-chain state, all parties can have a consistent view of the current active storage nodes.
Byzantine storage nodes can also freely join and leave the system.
However, we assume Byzantine entities are bounded by $\frac{1}{3}$ of the overall tokens, so at most $\frac{1}{3}$ of the nodes in the $\st{stakingSet}$ at any time can be Byzantine.
\lijl{Currently, we don't specify how much stake is required, and how that relates to the storage capacity of the node, or the rewards. We should discuss briefly.}

\subsubsection{Store Procedure}
\label{sec:design:proto:store}
Something.

\subsubsection{Data Retrieval}
\label{sec:design:proto:retrieval}
Something.

\subsubsection{Repair}
\label{sec:design:proto:repair}
Something.

\subsection{System Parameter Configuration}
\label{sec:design:param}
Something.

\subsection{Data Availability Proof}
\label{sec:design:proof}
Something.
